\let\R\relax
\newcommand{\R}{\mathbb{R}}
\let\C\relax
\newcommand{\C}{\mathbb{C}}
\let\N\relax
\newcommand{\N}{\mathbb{N}}
\let\Z\relax
\newcommand{\Z}{\mathbb{Z}}
\let\Q\relax
\newcommand{\Q}{\mathbb{Q}}
\let\SO\relax
\newcommand{\SO}{\operatorname{SO}}
\let\SL\relax
\newcommand{\SL}{\operatorname{SL}}
\let\GL\relax
\newcommand{\GL}{\operatorname{GL}}
\let\St\relax
\newcommand{\St}{\operatorname{St}}
\let\GV\relax
\newcommand{\GV}{\operatorname{GV}}
\let\E\relax
\newcommand{\E}{\operatorname{E}}
\let\Sp\relax
\newcommand{\Sp}{\operatorname{Sp}}
\let\Aut\relax
\newcommand{\Aut}{\operatorname{Aut}}
\let\ad\relax
\newcommand{\ad}{\operatorname{ad}}
\let\PSL\relax
\newcommand{\PSL}{\operatorname{PSL}}
\let\PSp\relax
\newcommand{\PSp}{\operatorname{PSp}}
\let\PSU\relax
\newcommand{\PSU}{\operatorname{PSU}}

\chapter{Jacques Tits (1993)}
\index{Tits, Jacques}

\begin{quote}
\it ``Tits has revolutionized group theory by introducing geometric ideas into the subject. His theory of buildings provides a unified geometric framework for understanding all groups of Lie type, and his classification of spherical buildings is one of the crowning achievements of 20th-century mathematics. The Tits alternative, a fundamental result about linear groups, has become a cornerstone of geometric group theory.'' --- Wolf Prize Citation
\end{quote}

Jacques Tits (1930--2021) was one of the most original mathematicians of the 20th century. The 1993 Wolf Prize in Mathematics (shared with Mikhail Gromov) recognized his fundamental contributions to group theory, geometry, and algebra. Tits created the theory of buildings---a geometric structure that generalizes projective spaces and provides a unified framework for understanding all groups of Lie type. He also proved the Tits alternative, a foundational result about the structure of linear groups, and developed the theory of BN-pairs, which axiomatically captures the structure of Lie-type groups. In 2008, he was awarded the Abel Prize for his lifetime achievements.

Tits's mathematical style is characterized by deep geometric insight combined with algebraic rigor. He had a remarkable ability to find the right level of abstraction to illuminate the structure of complicated algebraic objects. His work transformed the study of algebraic groups, Coxeter groups, and geometry, creating the modern field of incidence geometry and laying the foundations for what is now called geometric group theory.

\section{Early Life and Career}

Jacques Tits was born on August 12, 1930, in Uccle, a suburb of Brussels, Belgium. His father was a high school mathematics teacher. Tits showed extraordinary mathematical talent from an early age. He attended the Lyc\'ee in Brussels and then entered the Universit\'e Libre de Bruxelles (ULB) at age 15.

Tits completed his PhD in 1950 at ULB under the supervision of Alfred Errera. His doctoral thesis concerned the structure of algebraic groups and their geometries. He was only 20 years old. In 1951, he published his first major paper on the structure of algebraic groups, which already contained the seeds of his later theory of buildings.

From 1950 to 1962, Tits held positions at ULB, first as a researcher at the Belgian National Fund for Scientific Research and then as a professor. During this period, he developed the key ideas that would become the theory of buildings and BN-pairs. In 1962, he published his influential work on algebraic groups and their geometries, laying the groundwork for his systematic theory.

In 1962, Tits moved to the University of Bonn, where he remained until 1974. His time in Bonn was extraordinarily productive. He developed the axiomatic theory of buildings, classified spherical buildings, and proved the Tits alternative. In 1974, he moved to the Coll\`ege de France in Paris, where he held the chair of Group Theory until his retirement in 1995.

At the Coll\`ege de France, Tits shifted his focus to the construction of Kac--Moody groups and the development of the theory of buildings for arbitrary Coxeter groups. He remained active in research until late in life. He was awarded the Wolf Prize in 1993, the Abel Prize in 2008, and numerous other honors. He died on December 5, 2021, at the age of 91.

\section{BN-Pairs: The Axiomatic Structure of Lie-Type Groups}

One of Tits's most influential contributions is the concept of a BN-pair, also called a Tits system. This is an axiomatic framework that captures the essential structure of groups of Lie type, such as \(\SL_n(F)\), \(\SO_n(F)\), and the exceptional algebraic groups.

\begin{definition}[BN-Pair]
Let \(G\) be a group. A \textbf{BN-pair} (or \textbf{Tits system}) in \(G\) consists of two subgroups \(B\) and \(N\) of \(G\) satisfying the following axioms:
\begin{enumerate}
    \item \(G = \langle B, N \rangle\) and \(H = B \cap N\) is a normal subgroup of \(N\).
    \item The quotient \(W = N/H\) is generated by a set \(S = \{s_1, \dots, s_r\}\) of involutions. The pair \((W, S)\) is a \textbf{Coxeter system}.
    \item For every \(s \in S\) and every \(w \in W\),
    \[
    B s B \cdot B w B \subseteq B w B \cup B s w B.
    \]
    \item For every \(s \in S\), we have \(s B s \neq B\).
\end{enumerate}
\end{definition}

The subgroup \(B\) is called the \textbf{Borel subgroup}, \(N\) is the \textbf{monomial subgroup}, \(H = B \cap N\) is the \textbf{Cartan subgroup}, and \(W\) is the \textbf{Weyl group} of the Tits system.

\begin{example}[The Prototype: \(\SL_n(F)\)]
Let \(F\) be a field and let \(G = \SL_n(F)\). Let \(B\) be the subgroup of upper-triangular matrices in \(G\) (the Borel subgroup). Let \(N\) be the subgroup of monomial matrices (matrices with exactly one nonzero entry in each row and column, whose determinant is 1). Then \(H = B \cap N\) is the subgroup of diagonal matrices in \(\SL_n(F)\). The Weyl group \(W = N/H\) is isomorphic to the symmetric group \(S_n\), generated by the simple transpositions \(s_i = (i, i+1)\). The Coxeter presentation of \(S_n\) is:
\[
W = \langle s_1, \dots, s_{n-1} \mid s_i^2 = 1, (s_i s_{i+1})^3 = 1, (s_i s_j)^2 = 1 \text{ for } |i-j| \geq 2 \rangle.
\]
\end{example}

\begin{example}[The Symplectic Group \(\Sp_4(F)\)]
Let \(F\) be a field and let \(G = \Sp_4(F)\) be the group of \(4 \times 4\) matrices preserving the symplectic form \(\omega(x, y) = x_1 y_3 + x_2 y_4 - x_3 y_1 - x_4 y_2\). Let \(B\) be the subgroup of upper-triangular matrices in \(\Sp_4(F)\). Let \(N\) be the subgroup of monomial symplectic matrices. The Weyl group \(W = N/H\) is the dihedral group \(D_4\) of order 8, generated by two involutions \(s_1, s_2\) satisfying \((s_1 s_2)^4 = 1\). The building of \(\Sp_4(F)\) is a generalized 4-gon, whose apartments are regular octagons. This example illustrates how the Coxeter diagram (of type \(B_2 = C_2\)) determines the geometry of the building.
\end{example}

The axioms of a BN-pair lead to the \textbf{Bruhat decomposition}, one of the most important structural results in the theory of algebraic groups.

\begin{theorem}[Bruhat Decomposition]
Let \((B, N)\) be a BN-pair in \(G\) with Weyl group \(W\). Then
\[
G = \bigsqcup_{w \in W} B w B,
\]
a disjoint union of double cosets. Moreover, the multiplication of double cosets satisfies
\[
(B w B) \cdot (B w' B) \subseteq \bigcup_{v \leq w w'} B v B,
\]
where \(\leq\) denotes the Bruhat order on \(W\).
\end{theorem}

The Bruhat decomposition generalizes the classical decomposition of matrices into products of upper-triangular matrices and permutation matrices. It provides the fundamental combinatorial framework for understanding the structure of Lie-type groups.

The BN-pair axioms also give rise to the notion of \textbf{parabolic subgroups}. A subgroup \(P\) of \(G\) is called \textbf{parabolic} if it contains a conjugate of \(B\). Every parabolic subgroup corresponds to a subset \(J \subseteq S\) of the generating set of \(W\). The \textbf{Langlands decomposition} \(P = M_J N_J\) of a parabolic subgroup relates parabolic subgroups to Levi subgroups, which have their own BN-pair structure.

The existence of a BN-pair has profound consequences for the representation theory of \(G\). The \textbf{Hecke algebra} \(\mathcal{H}(G, B)\) of \(B\)-bi-invariant functions on \(G\) has a presentation in terms of the Weyl group and its Coxeter relations. The representations of \(G\) can be studied through the \textbf{principal series representations} induced from characters of \(B\), and the classification of irreducible representations reduces to the study of the Weyl group via the Kazhdan--Lusztig theory.

\section{Tits Buildings}

The theory of buildings, Tits's most famous creation, grew out of his work on algebraic groups and their geometries. A building is a combinatorial-geometric structure that generalizes projective spaces and provides a unified geometric framework for understanding all groups of Lie type.

\begin{definition}[Chamber System]
Let \(I\) be a finite set called the \textbf{type set}. A \textbf{chamber system} over \(I\) is a set \(\mathcal{C}\) of \textbf{chambers} together with a family of equivalence relations \(\sim_i\) for each \(i \in I\), called \(i\)-adjacency. Two chambers are \(i\)-adjacent if they are related by \(\sim_i\). A \textbf{gallery} is a finite sequence of chambers where consecutive chambers are \(i\)-adjacent for some \(i\).
\end{definition}

Buildings are chamber systems satisfying additional axioms that encode the structure of a Coxeter group.

\begin{definition}[Building]
Let \((W, S)\) be a Coxeter system with length function \(\ell\). A \textbf{building of type \((W, S)\)} is a chamber system \(\Delta\) over the set \(S\) equipped with a \textbf{Weyl distance function} \(\delta: \Delta \times \Delta \to W\) satisfying:
\begin{enumerate}
    \item \(\delta(C, D) = 1\) if and only if \(C = D\).
    \item If \(\delta(C, D) = w\) and \(C'\) is \(s\)-adjacent to \(C\) with \(C' \neq C\), then \(\delta(C', D) = s w\) or \(\delta(C', D) = w\). If \(\ell(s w) > \ell(w)\), then \(\delta(C', D) = s w\); if \(\ell(s w) < \ell(w)\), then \(\delta(C', D) = w\).
    \item For any chamber \(C\) and any \(w \in W\), there exists a chamber \(D\) such that \(\delta(C, D) = w\).
\end{enumerate}
\end{definition}

Equivalently, a building can be defined in terms of \textbf{apartments}. An apartment is a subcomplex isomorphic to the Coxeter complex of \((W, S)\). The building axioms require that any two chambers lie in a common apartment and that any two apartments are isomorphic via a unique isomorphism fixing their intersection.

\begin{definition}[Apartment Definition]
A \textbf{building} is a simplicial complex \(\Delta\) that can be expressed as a union of subcomplexes called \textbf{apartments}, each isomorphic to the Coxeter complex \(\Sigma(W, S)\), such that:
\begin{enumerate}
    \item Any two simplices of \(\Delta\) are contained in a common apartment.
    \item If \(A\) and \(A'\) are two apartments containing two simplices \(x\) and \(y\), then there is an isomorphism \(A \to A'\) fixing \(x\) and \(y\) pointwise.
\end{enumerate}
\end{definition}

\begin{example}[Projective Spaces as Buildings]
Let \(V\) be an \((n+1)\)-dimensional vector space over a field \(F\). The projective space \(\mathbb{P}^n(F)\) is a building of type \(A_n\). The chambers are the maximal flags
\[
V_1 \subset V_2 \subset \cdots \subset V_n \subset V,
\]
where \(\dim V_i = i\). Two chambers are \(i\)-adjacent if they differ only in the \(i\)-dimensional subspace. The Weyl group is the symmetric group \(S_{n+1}\). The apartments correspond to the choice of a basis for \(V\); each basis determines a set of coordinate subspaces, and these form a subcomplex isomorphic to the Coxeter complex of \(S_{n+1}\).
\end{example}

\begin{example}[Generalized \(m\)-gons]
A generalized \(m\)-gon is a building of rank 2 whose Weyl group is the dihedral group \(D_m\) of order \(2m\). Generalized \(m\)-gons are bipartite graphs of diameter \(m\) and girth \(2m\) with no circuits of length less than \(2m\). The simplest example is the ordinary \(m\)-gon, which is a building of type \(\tilde{A}_1\). Projective planes are generalized 3-gons. Classical examples arise from the groups \(\SL_3(F)\) (generalized 3-gons), \(\Sp_4(F)\) (generalized 4-gons), and \(G_2(F)\) (generalized 6-gons).
\end{example}

\subsection{Thick Buildings}

A building is called \textbf{thick} if every panel (a codimension-1 face of a chamber) is contained in at least three chambers. Thickness is the geometric analog of the condition that the group \(G\) acts transitively on chambers with a Borel subgroup \(B\) as the stabilizer of a chamber. In a thick building, the full automorphism group typically has a BN-pair.

The theory of thick buildings is central to Tits's classification. Every thick irreducible spherical building of rank at least 3 is associated to a semisimple algebraic group over a field, and its structure is completely determined by the group and the field.

\subsection{Spherical Buildings}

A building is called \textbf{spherical} if its Weyl group \(W\) is finite, equivalently, if the Coxeter complex \(\Sigma(W, S)\) is homeomorphic to a sphere. Spherical buildings arise naturally from algebraic groups over fields, from classical groups, and from exceptional groups.

\begin{theorem}[Tits Classification of Spherical Buildings, 1974]
Let \(\Delta\) be a thick irreducible spherical building of rank \(r \geq 3\). Then \(\Delta\) is isomorphic to the building associated to a semisimple algebraic group over a field, or to the building of a classical group over a division ring, or to one of the exceptional buildings associated to groups of type \(E_6, E_7, E_8, F_4\), or \(G_2\).

More precisely, the building \(\Delta\) is the flag complex of subspaces of a vector space, a polar space, or an exceptional geometry determined by the Dynkin diagram of the corresponding algebraic group.
\end{theorem}

This classification is one of the crowning achievements of 20th-century geometry. It demonstrates that all sufficiently symmetric geometries (thick, spherical, rank at least 3) arise from algebraic groups. The classification proceeds in two steps:
\begin{enumerate}
    \item First, the building determines a \textbf{Moufang polygon} or \textbf{Moufang building}, which satisfies a strong symmetry property (the Moufang condition).
    \item Second, the Moufang condition forces the building to be associated to an algebraic group by the work of Tits and Weiss.
\end{enumerate}

The classification for rank 2 (generalized polygons) is more intricate and was completed by Tits and Richard Weiss in their monograph \textit{Moufang Polygons}. The rank-2 case includes exotic examples that do not arise from algebraic groups over fields, such as those defined over alternative division algebras (octonions).

\subsection{The Moufang Property}

A building satisfies the \textbf{Moufang property} if for every root (a half-apartment), there is a subgroup of automorphisms fixing the root pointwise and acting transitively on the apartments containing that root. This property is a powerful symmetry condition that characterizes buildings arising from algebraic groups.

\begin{theorem}[Tits, 1974]
Every thick irreducible spherical building of rank at least 2 that satisfies the Moufang property is associated to a semisimple algebraic group over a field, a classical group over a division ring, or one of the exceptional groups.
\end{theorem}

The Moufang property is automatic for buildings of rank at least 3 by work of Tits, and for many buildings of rank 2. The classification of Moufang polygons by Tits and Weiss is a landmark in incidence geometry.

\section{The Tits Alternative}

The Tits alternative is a fundamental result in geometric group theory that describes the dichotomy in the structure of linear groups.

\begin{theorem}[Tits Alternative, 1972]
Let \(G\) be a finitely generated subgroup of \(\GL_n(K)\), where \(K\) is a field. Then either
\begin{enumerate}
    \item \(G\) contains a non-abelian free subgroup \(F_2\) (the free group on two generators), or
    \item \(G\) is \textbf{virtually solvable}: it contains a solvable subgroup of finite index.
\end{enumerate}
\end{theorem}

This dichotomy is striking: every linear group either is "large" in the sense of containing a free subgroup, or is "small" in the sense of being virtually solvable. The two possibilities are mutually exclusive: a solvable group cannot contain a non-abelian free subgroup, and a group containing a non-abelian free subgroup grows exponentially.

\subsection{Sketch of the Proof}

Tits's proof uses the dynamics of group actions on projective space \(\mathbb{P}^{n-1}(K)\). The key idea is to find two matrices \(a, b \in G\) that act as \textbf{proximal} elements: they have a unique attracting fixed point in \(\mathbb{P}^{n-1}(K)\) and a unique repelling fixed point, with all other points attracted toward the attracting fixed point under iteration.

The proof proceeds in several steps:
\begin{enumerate}
    \item If \(G\) is not virtually solvable, then its Zariski closure in \(\GL_n(K)\) is a non-solvable algebraic group.
    \item By a theorem of Borel, such a group contains a non-abelian free subgroup if and only if it contains a pair of elements generating a free subgroup.
    \item Using the \textbf{ping-pong lemma}, Tits constructed two matrices \(a, b \in G\) such that the subgroups they generate act on \(\mathbb{P}^{n-1}(K)\) with disjoint attracting and repelling sets, guaranteeing that they generate a free group.
    \item The ping-pong lemma requires the existence of elements with eigenvalues of different absolute values (in a suitable local field), which is ensured by the non-solvability of the Zariski closure.
\end{enumerate}

The ping-pong lemma, which Tits used in a masterful way, states: suppose a group \(\Gamma\) acts on a set \(X\) and there are elements \(\gamma_1, \gamma_2 \in \Gamma\) and nonempty disjoint subsets \(A_1, A_2 \subset X\) such that \(\gamma_i^{\pm 1}(X \setminus A_j) \subseteq A_i\) for \(i \neq j\). Then \(\gamma_1\) and \(\gamma_2\) generate a free subgroup of \(\Gamma\). Tits found such subsets in projective space by analyzing the dynamics of matrices with eigenvalues of distinct absolute values.

The proof also uses the \textbf{amenability} concept: virtually solvable groups are amenable, while groups containing a non-abelian free subgroup are not. This gives another perspective on the Tits alternative: a linear group is amenable if and only if it is virtually solvable. This characterization has important consequences in harmonic analysis and the theory of group actions.

\subsection{Consequences and Generalizations}

The Tits alternative has profound consequences across group theory:
\begin{itemize}
    \item It implies that any finitely generated linear group that is not virtually solvable has \textbf{exponential growth}.
    \item It gives a criterion for a linear group to be amenable: a finitely generated linear group is amenable if and only if it is virtually solvable.
    \item It shows that linear groups satisfy the \textbf{sofic property}: every linear group is sofic (a consequence of being residually finite, itself a property of linear groups).
\end{itemize}

The Tits alternative has been generalized to many other contexts:
\begin{itemize}
    \item \textbf{Gromov's alternative} for groups acting on hyperbolic spaces.
    \item The \textbf{Tits alternative for mapping class groups} (Ivanov, McCarthy): a subgroup of the mapping class group of a surface either contains a non-abelian free subgroup or is virtually abelian.
    \item The \textbf{Tits alternative for \(\text{Out}(F_n)\)} (Bestvina, Feighn, Handel): a subgroup of \(\text{Out}(F_n)\) either contains a non-abelian free subgroup or is virtually solvable.
    \item The \textbf{Tits alternative for CAT(0) groups} and groups acting on cube complexes.
\end{itemize}

The Tits alternative is now recognized as a fundamental property that characterizes "non-positively curved" groups. It is closely related to the concept of \textbf{rank} in group theory: groups of higher rank tend to contain free subgroups, while groups of rank 1 (such as hyperbolic groups) may or may not contain free subgroups depending on torsion.

\section{Buildings as Geometric Structures for Groups}

Tits's theory of buildings provides a geometric framework for studying groups. The key idea is that a group with a BN-pair acts transitively on a building, and the building encodes the combinatorial geometry of the group.

\subsection{The Building of a BN-Pair}

If \((B, N)\) is a BN-pair in a group \(G\) with Weyl group \((W, S)\), then the cosets of the Borel subgroup \(B\) form the chambers of a building. The building \(\Delta(G, B)\) is defined as follows:
\begin{itemize}
    \item The set of chambers is \(G/B\), the coset space of \(B\).
    \item Two chambers \(gB\) and \(hB\) are \(s\)-adjacent for \(s \in S\) if \(h^{-1}g \in B s B\).
    \item The Weyl distance is given by \(\delta(gB, hB) = w\) if and only if \(g^{-1}h \in B w B\) (the Bruhat decomposition).
\end{itemize}
The apartments are the subsets \(g N H / H \cong W\) for \(g \in G\), each isomorphic to the Coxeter complex of \((W, S)\).

This construction shows that every group with a BN-pair acts on a building by left multiplication, and the building is a geometric realization of the group's Bruhat decomposition.

\subsection{Group Actions on Buildings}

The theory of group actions on buildings is a rich subject. A group acting on a building by automorphisms can be studied through its action on the set of chambers and on apartments. Key results include:
\begin{itemize}
    \item \textbf{Bruhat--Tits fixed point theorem}: A group of bounded automorphisms of a building of non-positive curvature type has a global fixed point.
    \item \textbf{Simplicity criteria}: Many groups acting on buildings are simple, or have simple commutator subgroups. For example, if \(G\) is the group of automorphisms of a thick building that is topologically simple, then \(G\) is often simple.
    \item \textbf{The building at infinity}: For Euclidean buildings, the boundary at infinity is a spherical building, and the structure of this boundary encodes asymptotic properties of the group action.
\end{itemize}

\subsection{Buildings and Non-Positive Curvature}

Spherical buildings are examples of \textbf{non-positively curved} metric spaces. The apartments in a spherical building are isometric to the Coxeter complex \(\Sigma(W, S)\), which can be realized as a piecewise-spherical metric space. The building \(\Delta\) inherits a metric from its apartments, and this metric satisfies the \textbf{CAT(1)} condition: geodesic triangles in \(\Delta\) are at least as thin as their comparison triangles on the unit sphere.

Similarly, Euclidean buildings (Bruhat--Tits buildings) are \textbf{CAT(0)} spaces, meaning they have non-positive curvature in the sense of Alexandrov. The Bruhat--Tits tree for \(\SL_2(\Q_p)\) is a simple example: it is a tree, and all trees are CAT(0). Higher-rank Euclidean buildings are higher-dimensional analogues, with apartments isometric to Euclidean space.

The CAT(0) property of buildings has important consequences:
\begin{itemize}
    \item \textbf{Fixed point theorems}: If a group acts on a CAT(0) space by isometries and has a bounded orbit, then it has a global fixed point. This is the Bruhat--Tits fixed point theorem, which is essential for understanding the structure of \(p\)-adic groups.
    \item \textbf{Flat torus theorem}: Any abelian subgroup of a group acting properly on a CAT(0) space contains a free abelian subgroup of rank equal to the dimension of a flat in the space.
    \item \textbf{Rigidity theorems}: The CAT(0) structure of buildings is used in Mostow rigidity and in Margulis's superrigidity theorem for lattices in higher-rank groups.
\end{itemize}

\subsection{Applications to Algebraic Groups}

The theory of buildings was instrumental in Tits's work on algebraic groups. It provided a geometric language for studying parabolic subgroups, Levi decompositions, and the structure theory of reductive groups. The building of a reductive algebraic group \(G\) over a field \(F\) is called the \textbf{Tits building} of \(G(F)\).

For a reductive group \(G\), the Tits building \(\Delta(G(F))\) encodes the spherical Tits building associated to \(G\). The vertices of the building correspond to maximal parabolic subgroups, and the chambers correspond to Borel subgroups. The structure of the building reflects the Dynkin diagram of \(G\) and the field \(F\).

\subsection{Buildings and the Congruence Subgroup Problem}

Tits buildings also play a role in the congruence subgroup problem. The \textbf{Bruhat--Tits tree} for \(\SL_2\) over a local field is a tree (a rank-1 building) whose vertices correspond to lattices modulo homothety. The group \(\SL_2(F)\) acts on this tree, and the structure of the tree reveals the structure of the group's subgroups. The theory of buildings for higher-rank groups generalizes this to more complicated simplicial complexes, known as \textbf{Bruhat--Tits buildings} or \textbf{Euclidean buildings}.

\section{Kac--Moody Algebras and Groups}

In the 1980s, Tits turned his attention to the construction of groups associated to Kac--Moody Lie algebras. Kac--Moody algebras are infinite-dimensional Lie algebras that generalize finite-dimensional semisimple Lie algebras. They were discovered independently by Victor Kac and Robert Moody in the late 1960s.

\subsection{Kac--Moody Algebras}

A Kac--Moody algebra is defined by a generalized Cartan matrix \(A = (a_{ij})\) of size \(n \times n\) satisfying:
\begin{enumerate}
    \item \(a_{ii} = 2\) for all \(i\).
    \item \(a_{ij} \leq 0\) for \(i \neq j\).
    \item \(a_{ij} = 0\) if and only if \(a_{ji} = 0\).
\end{enumerate}
The associated Kac--Moody algebra \(\mathfrak{g}(A)\) is generated by \(3n\) generators \(e_i, f_i, h_i\) with relations:
\[
[h_i, h_j] = 0,\quad [h_i, e_j] = a_{ij} e_j,\quad [h_i, f_j] = -a_{ij} f_j,
\]
\[
[e_i, f_j] = \delta_{ij} h_i,
\]
\[
(\ad e_i)^{1-a_{ij}} e_j = 0,\quad (\ad f_i)^{1-a_{ij}} f_j = 0 \quad \text{for } i \neq j.
\]

When \(A\) is positive definite, \(\mathfrak{g}(A)\) is a finite-dimensional semisimple Lie algebra. When \(A\) is positive semidefinite, \(\mathfrak{g}(A)\) is an affine Kac--Moody algebra, which arises in the study of loop groups and integrable systems. When \(A\) is indefinite, \(\mathfrak{g}(A)\) is a more exotic infinite-dimensional algebra.

\subsection{Tits's Construction of Kac--Moody Groups}

While the algebraic theory of Kac--Moody algebras was well-developed by the 1970s, the corresponding group theory was less understood. Tits (1987) gave a general construction of Kac--Moody groups as groups with a BN-pair whose Weyl group is the Coxeter group associated to the generalized Cartan matrix.

Tits's construction proceeds by generators and relations. For each root \(\alpha\) of the Kac--Moody algebra, one defines a one-parameter subgroup \(U_\alpha\) of the group. The group is then generated by these root subgroups, subject to commutation relations that mirror the structure of the Kac--Moody algebra. The resulting group has a BN-pair \((B, N)\) where the Weyl group is the Coxeter group associated to the Cartan matrix.

The key insight is that Kac--Moody groups admit a geometric interpretation via buildings. The building of a Kac--Moody group is a \textbf{twin building}: a pair of buildings \((\Delta_+, \Delta_-)\) together with a codistance function \(\delta^*: (\Delta_+ \times \Delta_-) \cup (\Delta_- \times \Delta_+) \to W\) satisfying axioms that generalize the notion of a spherical building to the non-spherical case.

\begin{definition}[Twin Building]
A \textbf{twin building} of type \((W, S)\) is a pair \((\Delta_+, \Delta_-)\) of buildings of type \((W, S)\) together with a \textbf{codistance} function \(\delta^*: (\Delta_+ \times \Delta_-) \cup (\Delta_- \times \Delta_+) \to W\) such that:
\begin{enumerate}
    \item For all chambers \(C \in \Delta_\varepsilon\) and \(D \in \Delta_{-\varepsilon}\), we have \(\delta^*(C, D) = \delta^*(D, C)^{-1}\).
    \item If \(\delta^*(C, D) = w\) and \(C'\) is \(s\)-adjacent to \(C\) with \(C' \neq C\), then \(\delta^*(C', D) = s w\) or \(\delta^*(C', D) = w\). If \(\ell(s w) > \ell(w)\), then \(\delta^*(C', D) = s w\); if \(\ell(s w) < \ell(w)\), then \(\delta^*(C', D) = w\).
    \item For any \(C \in \Delta_\varepsilon\) and any \(w \in W\), there exists \(D \in \Delta_{-\varepsilon}\) such that \(\delta^*(C, D) = w\).
\end{enumerate}
\end{definition}

The existence of twin buildings for Kac--Moody groups was established by Tits and later developed by Markus Reineke, Bertrand R\'emy, and others. Twin buildings are central to the theory of \textbf{automorphic forms on Kac--Moody groups}, which generalizes the classical theory of automorphic forms on semisimple groups.

\subsection{Applications of Kac--Moody Groups}

Kac--Moody groups appear in many areas of mathematics:
\begin{itemize}
    \item \textbf{Mathematical physics}: Affine Kac--Moody groups are the symmetry groups of many integrable systems and conformal field theories. The representation theory of affine Kac--Moody algebras leads to the theory of vertex operator algebras and the Monster group (the moonshine conjecture).
    \item \textbf{Combinatorics}: The representation theory of Kac--Moody algebras involves the Kazhdan--Lusztig polynomials, which have deep connections to the combinatorics of Coxeter groups.
    \item \textbf{Geometric group theory}: Non-affine Kac--Moody groups are examples of groups with exotic geometric properties: they are finitely generated (when the Cartan matrix is of finite type) but have infinite-dimensional asymptotic cones and exhibit phenomena not seen in finite-dimensional groups.
\end{itemize}

\section{The Theory of Thick Buildings and Moufang Polygons}

The classification of thick buildings is one of Tits's most significant achievements. In rank 2, the buildings are \textbf{generalized polygons}: bipartite graphs with certain regularity conditions. The classification of Moufang polygons (generalized polygons satisfying the Moufang condition) was completed by Tits and Weiss.

\begin{theorem}[Tits--Weiss Classification of Moufang Polygons]
Every Moufang polygon arises from one of the following algebraic structures:
\begin{itemize}
    \item A field, for type \(A_2\) (projective planes) and \(B_2\) (quadrangles of orthogonal type).
    \item A quadratic space, for type \(B_2\) (symplectic quadrangles) and \(C_2\).
    \item A Cayley division algebra (octonions), for type \(G_2\) (hexagons) and \(F_4\) (octagons).
    \item An anisotropic form of type \(E_6, E_7, E_8\), yielding exceptional polygons.
\end{itemize}
\end{theorem}

The Moufang condition is equivalent to the existence of a group of automorphisms that acts transitively on the set of apartments containing a given root. This condition is automatically satisfied for thick buildings of rank at least 3, but for rank 2 it is an additional hypothesis that picks out the most symmetric polygons.

The classification of Moufang polygons is a monumental achievement that required a detailed analysis of the possible root group structures. The result is a complete list of possibilities, each corresponding to a well-defined algebraic object. This classification, together with the classification of spherical buildings of higher rank, gives a complete description of all highly symmetric incidence geometries.

\section{Impact and Legacy}

\subsection{The Theory of Buildings}

Tits's theory of buildings has become an essential tool in many areas of mathematics:
\begin{itemize}
    \item \textbf{Algebraic groups}: The building of a reductive algebraic group encodes its parabolic subgroups and their interrelations. The Bruhat--Tits theory of groups over local fields uses Euclidean buildings to study the structure of \(p\)-adic groups.
    \item \textbf{Representation theory}: The geometry of buildings is used in the study of discrete series representations, the construction of Hecke algebras, and the Kazhdan--Lusztig theory.
    \item \textbf{Rigidity theory}: The classification of spherical buildings is used in Mostow rigidity, superrigidity theorems (Margulis), and the study of lattices in semisimple groups.
    \item \textbf{Geometric group theory}: Buildings provide some of the most important examples of CAT(0) spaces and spaces with non-positive curvature.
\end{itemize}

\subsection{The Tits Alternative}

The Tits alternative is a cornerstone of geometric group theory. It has been generalized to mapping class groups, \(\text{Out}(F_n)\), groups acting on hyperbolic spaces, and many other contexts. It provides a fundamental dichotomy that organizes the structure of large classes of groups.

\subsection{BN-Pairs}

The concept of a BN-pair has become a standard tool in group theory. It provides the axiomatic foundation for the theory of finite groups of Lie type (the classification of finite simple groups relies heavily on BN-pair theory). The Bruhat decomposition, parabolic subgroups, and the theory of Hecke algebras all flow from the BN-pair axioms.

\subsection{Bruhat--Tits Theory}

The theory of Bruhat--Tits buildings (Euclidean buildings) for groups over local fields is a major chapter in number theory and representation theory. It provides the geometric foundation for the study of \(p\)-adic groups, their representations, and their cohomology. The \textbf{Bruhat--Tits tree} for \(\SL_2(\Q_p)\) is a simple but powerful tool for understanding the structure of the group \(\SL_2(\Q_p)\) and its subgroups.

\subsection{Education and Exposition}

Tits was a master of clear exposition. His writing style is precise and elegant, and his books and papers have educated generations of mathematicians. His monograph \textit{Lectures on Buildings} (with F. Bruhat) remains a standard reference, and his papers on the classification of spherical buildings are models of clarity and rigor.

Tits had a profound influence on his students and collaborators. Among his most famous students are Philippe Gille, Bertrand R\'emy, and Richard Weiss. His collaboration with Richard Weiss on the classification of Moufang polygons produced one of the deepest results in incidence geometry.

\subsection{Tits and the Classification of Finite Simple Groups}

The theory of BN-pairs played a crucial role in the classification of finite simple groups, one of the greatest achievements of 20th-century mathematics. Every finite group of Lie type admits a BN-pair, and the classification of finite simple groups uses the BN-pair structure to identify groups of Lie type among the finite simple groups. The parameters of the BN-pair (the rank, the Weyl group, and the field size) distinguish the different families: \(\PSL_n(q)\), \(\PSp_{2n}(q)\), \(\PSU_n(q)\), the exceptional groups \(E_6(q), E_7(q), E_8(q), F_4(q), G_2(q)\), and the twisted groups \(^2E_6(q), ^3D_4(q), ^2B_2(q), ^2G_2(q), ^2F_4(q)\). In each case, the associated building determines the group uniquely.

\subsection{Continuing Influence}

Tits's ideas continue to shape mathematical research:
\begin{itemize}
    \item The theory of buildings is being extended to \textbf{universal buildings} and \textbf{timber buildings}, which generalize the classical theory to new contexts.
    \item The Tits alternative continues to inspire new results in the study of group actions on hyperbolic and CAT(0) spaces.
    \item Kac--Moody groups are an active area of research, with connections to automorphic forms, mathematical physics, and geometric group theory.
    \item The theory of twin buildings is being developed as a tool for understanding the geometry of Kac--Moody groups and their representations.
\end{itemize}

\section{Timeline}

\begin{tabularx}{\linewidth}{|l|X|}
\hline
\textbf{Year} & \textbf{Event} \\ \hline
1930 & Born on August 12 in Uccle, Brussels, Belgium \\ \hline
1945 & Entered the Universit\'e Libre de Bruxelles at age 15 \\ \hline
1950 & Ph.D. in Mathematics, Universit\'e Libre de Bruxelles (advisor: Alfred Errera) \\ \hline
1950 & Published first major paper on algebraic groups and their geometries \\ \hline
1951--1962 & Researcher and Professor at ULB; developed early ideas on buildings \\ \hline
1962 & Moved to the University of Bonn as Professor of Mathematics \\ \hline
1962 & Published foundational papers on BN-pairs and the Bruhat decomposition \\ \hline
1969 & Published \textit{Lectures on Buildings}, systematizing the theory \\ \hline
1972 & Published the Tits alternative for linear groups \\ \hline
1974 & Moved to the Coll\`ege de France, Paris (Chair of Group Theory) \\ \hline
1974 & Completed the classification of spherical buildings \\ \hline
1980s & Developed the theory of Kac--Moody groups and twin buildings \\ \hline
1987 & Published the construction of Kac--Moody groups \\ \hline
1993 & Awarded the Wolf Prize in Mathematics (shared with Mikhail Gromov) \\ \hline
1995 & Retired from the Coll\`ege de France \\ \hline
2008 & Awarded the Abel Prize from the Norwegian Academy of Science and Letters \\ \hline
2021 & Died on December 5 at age 91 \\ \hline
\end{tabularx}

\section{Frequently Asked Questions}

\subsection*{Q1: What is a Tits building in intuitive terms?}

A Tits building is a geometric structure that generalizes projective spaces. Think of it as a giant "apartment building": the building is made up of many apartments (each isomorphic to a Coxeter complex), and any two rooms (chambers) are located in some common apartment. The apartments fit together in a highly symmetric way. Buildings provide a geometric language for studying groups with BN-pairs: the group acts transitively on the chambers, and the action preserves the building structure.

\subsection*{Q2: What is the Tits alternative saying about linear groups?}

The Tits alternative says that any finitely generated linear group is either "big" (contains a non-abelian free group) or "small" (virtually solvable). This is a stark dichotomy: there is no intermediate possibility. For example, \(\SL_2(\Z)\) contains a free subgroup (it is big), while the Heisenberg group (upper triangular integer matrices with ones on the diagonal) is virtually solvable (it is small). The Tits alternative shows that the only way a linear group can avoid containing a free subgroup is by being essentially solvable.

\subsection*{Q3: What are BN-pairs and why are they important?}

A BN-pair (Tits system) is an axiomatic description of the structure of a group of Lie type. The two subgroups \(B\) (Borel) and \(N\) (monomial) generate the group, and their intersection \(H\) is the Cartan subgroup. The quotient \(W = N/H\) is the Weyl group, a Coxeter group. The axioms encode the Bruhat decomposition \(G = \bigsqcup_{w \in W} B w B\). BN-pairs are important because they provide a unified framework for studying all groups of Lie type, including finite simple groups. The classification of finite simple groups uses BN-pair theory in an essential way.

\subsection*{Q4: How does the classification of spherical buildings work?}

Tits classified all thick irreducible spherical buildings of rank at least 3. The classification says that every such building is the building of a semisimple algebraic group over a field, a classical group over a division ring, or one of the exceptional types \(E_6, E_7, E_8, F_4, G_2\). The proof proceeds by reconstructing the field and the group from the geometry of the building. The key idea is that the building determines a system of coordinates, and the Moufang condition (a symmetry property) forces these coordinates to come from an algebraic structure (a field, a division ring, or an alternative algebra).

\subsection*{Q5: What is a Kac--Moody group?}

A Kac--Moody group is a group associated to a Kac--Moody algebra, just as a semisimple algebraic group is associated to a finite-dimensional semisimple Lie algebra. Kac--Moody algebras generalize the finite-dimensional theory by allowing generalized Cartan matrices that are not positive definite. Tits constructed Kac--Moody groups using generators and relations, and showed that they have BN-pairs whose Weyl groups are infinite Coxeter groups. These groups act on twin buildings, which are pairs of buildings with a codistance function.

\subsection*{Q6: What are the major open problems related to Tits's work?}

\begin{enumerate}
    \item \textbf{Classification of rank-2 buildings:} While Moufang polygons are classified, not all thick generalized polygons are Moufang. The complete classification of all thick buildings of rank 2 (generalized \(m\)-gons) remains open for \(m \geq 3\).
    \item \textbf{Kac--Moody groups over finite fields:} The structure of Kac--Moody groups over finite fields and their representation theory is an active area of research with many open questions.
    \item \textbf{Buildings and the classification of groups:} The relationship between buildings, algebraic groups, and the classification of finite simple groups continues to be refined.
    \item \textbf{The congruence subgroup problem:} The Bruhat--Tits theory of buildings for groups over local fields is used in the study of the congruence subgroup problem, which remains open for many groups.
\end{enumerate}

\subsection*{Q7: How is the Tits alternative related to geometric group theory?}

The Tits alternative is one of the foundational results of geometric group theory. It shows that linear groups lie at one end of a spectrum: they either contain a free subgroup (and thus have exponential growth, are non-amenable, and have complicated Cayley graphs) or they are virtually solvable (and thus have polynomial growth, are amenable, and have relatively simple structure). The alternative has been generalized to many other contexts, and it serves as a model for the kind of structural dichotomies that geometric group theory seeks.

\section{Related Chapters}

See also Chapter~26 (Thompson) on finite simple groups.

\section*{Exercises}
\addcontentsline{toc}{section}{Exercises}

\begin{enumerate}
    \item [I] Show that in the BN-pair of \(\SL_n(F)\), the Weyl group is isomorphic to the symmetric group \(S_n\) and the Coxeter generators are the simple transpositions.
    \item [B] Prove that the Bruhat decomposition for \(\SL_2(F)\) takes the form \(\SL_2(F) = B \cup B w B\), where \(w = \begin{pmatrix} 0 & -1 \\ 1 & 0 \end{pmatrix}\).
    \item [I] Let \(\Delta\) be the building of \(\SL_3(F)\). Show that the apartments correspond to the choice of a basis of \(F^3\).
    \item [I] Prove that the projective plane \(\mathbb{P}^2(F)\) is a generalized 3-gon (building of type \(A_2\)).
    \item [B] Show that the Tits alternative holds in the trivial case: any finite group is virtually solvable (indeed, finite) and does not contain a non-abelian free subgroup.
    \item [I] Let \(G = \SL_2(\Z)\). Find an explicit non-abelian free subgroup of \(G\).
    \item [B] Prove that the free group \(F_2\) is not virtually solvable. Conclude that \(F_2\) satisfies the Tits alternative (it contains a non-abelian free subgroup, namely itself).
    \item [B] Show that the Heisenberg group \(H = \left\{ \begin{pmatrix} 1 & a & b \\ 0 & 1 & c \\ 0 & 0 & 1 \end{pmatrix} : a,b,c \in \Z \right\}\) is solvable and hence virtually solvable.
    \item [I] Let \(\Delta\) be a thick building of type \(A_n\). Show that \(\Delta\) is isomorphic to the flag complex of an \((n+1)\)-dimensional vector space over some division ring.
    \item [I] Show that the ordinary \(m\)-gon (an \(m\)-cycle graph) is a building of type \(\tilde{A}_1\) with the dihedral Weyl group \(D_m\).
    \item [I] Compute the Coxeter diagram for the Weyl group of type \(E_6\). Describe the corresponding building.
    \item [I] Show that the building of \(\Sp_4(F)\) is a generalized 4-gon. Describe its apartments.
    \item [I] Prove that a group with a BN-pair is transitive on the chambers of its building.
    \item [I] Let \(G\) be a group with a BN-pair. Show that the parabolic subgroups of \(G\) correspond to subsets of \(S\).
    \item [B] Show that the group of upper-triangular matrices in \(\SL_n(F)\) is the Borel subgroup \(B\) and that it is solvable.
    \item [I] Let \(g \in \GL_n(\C)\) be a matrix with eigenvalues of different absolute values. Show that the action of \(g\) on projective space \(\mathbb{P}^{n-1}(\C)\) has a unique attracting fixed point.
    \item [I] Prove the ping-pong lemma: if two elements \(a, b\) in a group acting on a set \(X\) have disjoint attracting and repelling sets, then they generate a free group.
    \item [I] Show that the building of a reductive algebraic group over a finite field is a finite simplicial complex.
    \item [I] Let \((W, S)\) be a Coxeter system. Show that the Coxeter complex \(\Sigma(W, S)\) is a building of type \((W, S)\).
    \item [I] Show that the Bruhat--Tits tree for \(\SL_2(\Q_p)\) is a regular tree of degree \(p+1\).
    \item [I] Let \(\mathfrak{g}(A)\) be an affine Kac--Moody algebra of type \(\tilde{A}_1\). Show that the Dynkin diagram of \(\mathfrak{g}(A)\) is an extended \(A_1\) diagram.
    \item [B] Show that a generalized \(m\)-gon has girth \(2m\) and diameter \(m\) as a bipartite graph.
    \item [I] Prove that in a building, any two chambers are connected by a gallery of minimal length equal to the length of the Weyl distance between them.
    \item [I] Let \(G\) be a semisimple algebraic group over a field \(F\). Describe the relationship between the Tits building of \(G(F)\) and the Dynkin diagram of \(G\).
    \item [B] Prove that the building of \(\SL_n(F)\) is thick if and only if \(F\) has at least three elements.
    \item [I] Show that the Weyl group of type \(G_2\) is the dihedral group \(D_6\) of order 12, and describe the corresponding generalized hexagon.
    \item [I] Let \(\Delta\) be a spherical building of type \(A_n\). Show that the apartments are isomorphic to the flag complex of an \((n+1)\)-dimensional vector space.
\end{enumerate}

\section*{Glossary}
\addcontentsline{toc}{section}{Glossary}

\begin{description}
    \item[Apartment] A subcomplex of a building isomorphic to the Coxeter complex of the Weyl group. Apartments are the "coordinate patches" of a building.
    \item[BN-pair (Tits system)] A pair of subgroups \((B, N)\) of a group satisfying axioms that capture the structure of groups of Lie type.
    \item[Borel subgroup] A maximal solvable subgroup of an algebraic group; in a BN-pair, the subgroup \(B\).
    \item[Bruhat decomposition] The decomposition \(G = \bigsqcup_{w \in W} B w B\) of a group with a BN-pair into double cosets indexed by the Weyl group.
    \item[Building] A combinatorial-geometric structure generalizing projective spaces, used to understand groups of Lie type.
    \item[Cartan subgroup] The intersection \(H = B \cap N\) in a BN-pair; a maximal torus in the algebraic group context.
    \item[Chamber] A maximal simplex in a building; the fundamental object of incidence geometry.
    \item[Coxeter complex] The simplicial complex associated to a Coxeter system \((W, S)\).
    \item[Coxeter system] A pair \((W, S)\) where \(W\) is generated by involutions \(s \in S\) with relations \((s_i s_j)^{m_{ij}} = 1\).
    \item[Generalized \(m\)-gon] A building of rank 2 with dihedral Weyl group \(D_m\); a bipartite graph of diameter \(m\) and girth \(2m\).
    \item[Kac--Moody algebra] An infinite-dimensional Lie algebra defined by a generalized Cartan matrix, generalizing semisimple Lie algebras.
    \item[Kac--Moody group] A group associated to a Kac--Moody algebra, having a BN-pair with an infinite Weyl group.
    \item[Moufang building] A building satisfying the Moufang property: for every root, a subgroup of automorphisms fixes the root pointwise and acts transitively on apartments.
    \item[Parabolic subgroup] A subgroup containing a conjugate of the Borel subgroup \(B\) in a BN-pair.
    \item[Ping-pong lemma] A criterion for two elements in a group action to generate a free subgroup, based on disjoint attracting and repelling sets.
    \item[Spherical building] A building whose Weyl group is finite; equivalently, the Coxeter complex is homeomorphic to a sphere.
    \item[Thick building] A building where every panel is contained in at least three chambers.
    \item[Tits alternative] The theorem that a finitely generated linear group either contains a non-abelian free subgroup or is virtually solvable.
    \item[Tits building] The building associated to an algebraic group, encoding its parabolic subgroups and their incidence geometry.
    \item[Twin building] A pair of buildings \((\Delta_+, \Delta_-)\) together with a codistance function, generalizing spherical buildings to non-spherical types.
    \item[Virtually solvable] A group having a solvable subgroup of finite index.
    \item[Weyl group] The quotient \(W = N/H\) in a BN-pair; a Coxeter group acting as the "reflection group" of the building.
\end{description}
